\documentclass[10pt, a4paper]{article}
\usepackage{setspace} 
\usepackage{amsmath, amssymb} 
\usepackage{geometry}
\usepackage{graphicx} 
\usepackage{slashed}
\usepackage{simplewick}
\doublespacing 

\title{Theory of Solid Note}
\author{Yun-Yen Wu} 
\date{} 

\begin{document}

\maketitle

\begin{abstract}
The "Theory of Solid" course is to study collective phenomena in condensed matter system and to solve a system with nuclei and electrons, taking into account their interactions (Coulomb). 
\end{abstract}

\newpage 

\setcounter{page}{1} 

\section{Introduction}

\section{Crystal Structure}

A perfect crystal is an array of atoms arranged in a \textbf{periodic} structure spanning the entire space. In this idealized model, we do not take surface physics, nanostructures, impurities, or disorders into consideration.

\subsection*{Fundamental Definitions}
The structure of all crystals can be described in terms of a lattice and a basis:
\begin{itemize}
    \item \textbf{Lattice:} An infinite and periodic array of discrete points in space.
    \item \textbf{Basis:} A group of atoms attached to each lattice point.
    \item \textbf{Crystal Structure = Lattice + Basis}
\end{itemize}

\subsection*{Real-Space Concepts}
\begin{description}
    \item[Bravais Lattice:] A $d$-dimensional Bravais lattice is an infinite array of discrete points generated by a set of $d$ linearly independent primitive vectors $\mathbf{a}_j$. The lattice consists of all position vectors $\mathbf{R}$ of the form:
    \begin{equation}
        \mathbf{R} = \sum^d_{j=1} n_j \mathbf{a}_j \quad \text{with}\quad n_j\in\mathbb{Z}
    \end{equation}
    A fundamental property is its \textbf{discrete translational symmetry}: the arrangement and orientation of the lattice must appear exactly the same when viewed from any arbitrary point. For a given lattice, the choice of primitive basis vectors is not unique.
    
    \item[Primitive Unit Cell:] The minimum volume of space that can fill all of space without overlapping when translated through all vectors in a Bravais lattice. It contains exactly one lattice point.
    
    \item[Wigner-Seitz Cell:] A special type of primitive cell defined as the region of space that is closer to a specific point of the Bravais lattice than to any other point.
    
    \item[Miller Indices:] Used to specify the orientation of a crystal plane, determined by three non-collinear points in the plane. A specific plane is denoted by parentheses $(hkl)$, while a family of planes equivalent by symmetry is denoted by curly brackets $\{hkl\}$.
\end{description}

\subsection*{Momentum Space (Reciprocal Space)}
\begin{description}
    \item[Reciprocal Lattice:] A set of all wave vectors $\mathbf{K}$ that yield plane waves (living in momentum space) with the periodicity of a given Bravais lattice $\mathbf{R}$, satisfying $e^{i\mathbf{R}\cdot\mathbf{K}}=1$. Essentially, it is the Fourier transform of the real-space crystal lattice and is itself a Bravais lattice.
    
    \item[First Brillouin Zone:] The Wigner-Seitz primitive cell of the reciprocal lattice. It is the fundamental volume in momentum space that contains all the unique wavevectors describing how quantum waves travel through the crystal.
\end{description}

\subsection*{Crystal Symmetries}
Symmetries of the crystal (e.g., translation, inversion, glide planes, screw axes) can be used to classify different crystals. A crystal's space group is determined by the set $\mathcal{G}$ of all symmetry operations that leave the body invariant.

\subsection*{Classification of Crystal Systems}

\begin{description}
    \item[Point Group:] All operations that transform the lattice into itself and leave a given point invariant. The point group of a lattice can not efine the lattice. 
    \item[Schonflies Notation:] 
    \begin{enumerate}
        \item $C_n$: Cyclic, groups having only one n-fold rotation axis
        \item $D_n$: Dihedral, group having n 2-fold rotation axia perpendicular to the pricipal $C_n$ axis
        \item $h$: horizontal, perpenicular to the rotation axis
        \item $v$: vertical, parallel to the rotation axis
        \item $d$: diagonal, parallel to the main rotation axis in the plane and bisecting the angle between the 2-fold axes perpendicular to the principal axis
        \item $S_n$: Spiegel, 
        \item $O$: Octahedral,
        \item $T$: Tetrahedral
    \end{enumerate}
\end{description}






\end{document}